% Options for packages loaded elsewhere
\PassOptionsToPackage{unicode}{hyperref}
\PassOptionsToPackage{hyphens}{url}
\PassOptionsToPackage{dvipsnames,svgnames,x11names}{xcolor}
\documentclass[
  10pt,
]{article}
\usepackage{xcolor}
\usepackage[margin=1cm,top=1cm,bottom=1cm,left=1cm,right=1cm,includeheadfoot]{geometry}
\usepackage{amsmath,amssymb}
\setcounter{secnumdepth}{3}
\usepackage{iftex}
\ifPDFTeX
  \usepackage[T1]{fontenc}
  \usepackage[utf8]{inputenc}
  \usepackage{textcomp} % provide euro and other symbols
\else % if luatex or xetex
  \usepackage{unicode-math} % this also loads fontspec
  \defaultfontfeatures{Scale=MatchLowercase}
  \defaultfontfeatures[\rmfamily]{Ligatures=TeX,Scale=1}
\fi
\usepackage{lmodern}
\ifPDFTeX\else
  % xetex/luatex font selection
  \setmainfont[]{DejaVu Serif}
  \setmonofont[]{DejaVu Sans Mono}
\fi
% Use upquote if available, for straight quotes in verbatim environments
\IfFileExists{upquote.sty}{\usepackage{upquote}}{}
\IfFileExists{microtype.sty}{% use microtype if available
  \usepackage[]{microtype}
  \UseMicrotypeSet[protrusion]{basicmath} % disable protrusion for tt fonts
}{}
\usepackage{setspace}
\makeatletter
\@ifundefined{KOMAClassName}{% if non-KOMA class
  \IfFileExists{parskip.sty}{%
    \usepackage{parskip}
  }{% else
    \setlength{\parindent}{0pt}
    \setlength{\parskip}{6pt plus 2pt minus 1pt}}
}{% if KOMA class
  \KOMAoptions{parskip=half}}
\makeatother
\usepackage{listings}
\newcommand{\passthrough}[1]{#1}
\lstset{defaultdialect=[5.3]Lua}
\lstset{defaultdialect=[x86masm]Assembler}
\usepackage{graphicx}
\makeatletter
\newsavebox\pandoc@box
\newcommand*\pandocbounded[1]{% scales image to fit in text height/width
  \sbox\pandoc@box{#1}%
  \Gscale@div\@tempa{\textheight}{\dimexpr\ht\pandoc@box+\dp\pandoc@box\relax}%
  \Gscale@div\@tempb{\linewidth}{\wd\pandoc@box}%
  \ifdim\@tempb\p@<\@tempa\p@\let\@tempa\@tempb\fi% select the smaller of both
  \ifdim\@tempa\p@<\p@\scalebox{\@tempa}{\usebox\pandoc@box}%
  \else\usebox{\pandoc@box}%
  \fi%
}
% Set default figure placement to htbp
\def\fps@figure{htbp}
\makeatother
\setlength{\emergencystretch}{3em} % prevent overfull lines
\providecommand{\tightlist}{%
  \setlength{\itemsep}{0pt}\setlength{\parskip}{0pt}}
% Enable graphics inclusion and ensure figure numbering works
\usepackage{graphicx}
\renewcommand{\figurename}{Figure}

% Configure fonts for Unicode support with fallbacks
\usepackage{newunicodechar}
\newunicodechar{⁴}{\textsuperscript{4}}
\newunicodechar{₄}{\textsubscript{4}}

% Math notation macros
\newcommand{\norm}[1]{\lVert #1\rVert}
\newcommand{\abs}[1]{\lvert #1\rvert}

% Enhanced code block styling for better contrast and readability
\usepackage{fancyvrb}
\usepackage{xcolor}
\usepackage{listings}

% Define custom colors for code blocks
\definecolor{codebg}{RGB}{245, 245, 245}      % Light gray background
\definecolor{codeborder}{RGB}{200, 200, 200}  % Medium gray border
\definecolor{codefg}{RGB}{50, 50, 50}         % Dark gray text

% Configure Verbatim environment for inline code
\DefineVerbatimEnvironment{Verbatim}{Verbatim}{%
    fontsize=\small,
    frame=single,
    framerule=0.5pt,
    framesep=3pt,
    rulecolor=\color{codeborder},
    bgcolor=\color{codebg},
    fgcolor=\color{codefg}
}

% Configure code block styling
\DefineVerbatimEnvironment{Highlighting}{Verbatim}{%
    fontsize=\footnotesize,
    frame=single,
    framerule=0.5pt,
    framesep=5pt,
    rulecolor=\color{codeborder},
    bgcolor=\color{codebg},
    fgcolor=\color{codefg}
}

% Style inline code with \texttt
\renewcommand{\texttt}[1]{%
    \colorbox{codebg}{\color{codefg}\ttfamily #1}%
}

% Configure listings package for code blocks
\lstset{
    backgroundcolor=\color{codebg},
    basicstyle=\footnotesize\ttfamily\color{codefg},
    breakatwhitespace=false,
    breaklines=true,
    captionpos=b,
    commentstyle=\color{codefg},
    deletekeywords={...},
    escapeinside={\%*}{*)},
    extendedchars=true,
    frame=single,
    framerule=0.5pt,
    framesep=5pt,
    keepspaces=true,
    keywordstyle=\color{codefg},
    language=Python,
    morekeywords={*,...},
    numbers=left,
    numbersep=5pt,
    numberstyle=\tiny\color{codefg},
    rulecolor=\color{codeborder},
    showspaces=false,
    showstringspaces=false,
    showtabs=false,
    stepnumber=1,
    stringstyle=\color{codefg},
    tabsize=2,
    title=\lstname
}

% Override any Pandoc default lstset configurations
\AtBeginDocument{
    \lstset{
        backgroundcolor=\color{codebg},
        basicstyle=\footnotesize\ttfamily\color{codefg},
        frame=single,
        framerule=0.5pt,
        framesep=5pt,
        rulecolor=\color{codeborder},
        numbers=left,
        numbersep=5pt,
        numberstyle=\tiny\color{codefg}
    }
}

% Configure hyperref colors consistently
\AtBeginDocument{
% Override pandoc's hidelinks setting with consistent options
\hypersetup{
    colorlinks=true,
    allcolors=red,
    linkcolor=red,
    urlcolor=red,
    citecolor=red,
    filecolor=red,
    menucolor=red,
    linktoc=all
}
}

% Simple page break support for document structure
\usepackage{bookmark}
\IfFileExists{xurl.sty}{\usepackage{xurl}}{} % add URL line breaks if available
\urlstyle{same}
% fallback for those not using the hyperref driver hyperxmp:
\makeatletter
\@ifundefined{xmpquote}{\newcommand{\xmpquote}[1]{#1}}{}
\makeatother
\hypersetup{
  colorlinks=true,
  linkcolor={red},
  filecolor={red},
  citecolor={red},
  urlcolor={red},
  pdfcreator={LaTeX via pandoc}}

\title{Dynamics and Learning on the IVM Lattice}
\author{Daniel Ari Friedman\\ ORCID: 0000-0001-6232-9096\\ Email: daniel@activeinference.institute\\ DOI: 10.5281/zenodo.16887791}
\date{October 07, 2026}

\begin{document}
\maketitle

{
\hypersetup{linkcolor=black}
\setcounter{tocdepth}{3}
\tableofcontents
}
\setstretch{1.0}
\section{Dynamics and Learning on the IVM
Lattice}\label{dynamics-and-learning-on-the-ivm-lattice}

\subsection{Overview}\label{overview}

Where \href{04_optimization_in_4d.md}{Section 4} descends a \emph{single
point} along the IVM lattice, this section evolves a \emph{field} ---
one scalar per lattice site --- and then asks the inverse question:
given an observed field trajectory, what coupling produced it?
Everything here is implemented in
\passthrough{\lstinline!ivm\_dynamics.py!}
(\passthrough{\lstinline!ball\_sites!},
\passthrough{\lstinline!make\_lattice!},
\passthrough{\lstinline!heat\_step!},
\passthrough{\lstinline!majority\_step!},
\passthrough{\lstinline!step!}, \passthrough{\lstinline!simulate!},
\passthrough{\lstinline!fit\_trajectory!},
\passthrough{\lstinline!sum\_of\_squares!},
\passthrough{\lstinline!is\_nonincreasing!}) on top of the
\passthrough{\lstinline!Quadray!} class and
\passthrough{\lstinline!to\_xyz!} embedding in
\passthrough{\lstinline!quadray.py!}; the demo figure is produced by
\passthrough{\lstinline!quadmath/scripts/ivm\_dynamics\_demo.py!} (which
delegates to \passthrough{\lstinline!render\_dynamics\_demo!} in
\passthrough{\lstinline!src/quadmath/lattice/ivm\_dynamics.py!}, per the
thin-orchestrator contract in
\passthrough{\lstinline!quadmath/scripts/AGENTS.md!}).

\subsection{Lattice sites, shells, and the 12-around-one move
graph}\label{lattice-sites-shells-and-the-12-around-one-move-graph}

A \emph{site} is a canonical quadray representative --- non-negative
integer components with at least one zero, selected by
\passthrough{\lstinline!Quadray.normalize!} (see
\href{03_quadray_methods.md}{Quadray Methods}). The radial observable is
the squared embedding norm
\passthrough{\lstinline!site\_radius\_sq(q)!}; because the shift vector
\((1,1,1,1)\) lies in the kernel of the
\passthrough{\lstinline!DEFAULT\_EMBEDDING!} matrix, the radius is
invariant under quadray normalization.
\passthrough{\lstinline!ball\_sites(R)!} enumerates the canonical sites
with squared radius at most \(R^2\): for a canonical site the squared
radius is at least twice the square of its largest component, so
components in \([0, R]\) cover the ball.

\passthrough{\lstinline!make\_lattice(R)!} joins two sites when one of
the 12 canonical IVM neighbor shifts
(\passthrough{\lstinline!neighbor\_shifts!}, all permutations of
\((2,1,1,0)\), each at squared radius 8 --- the close-packing distance)
maps one site onto the other after renormalization. Normalization never
moves the embedded point, so adjacency is exactly ``embedding difference
is one of the 12 shifts''. Lattice sites are the quadrays whose
coordinate sum is divisible by 4; the remaining residue classes are IVM
voids and are excluded. On the \(R = 3\) ball (13 sites) the graph is a
single connected component: the close-packed \textbf{12-around-one
cluster}, i.e.~the origin (degree 12) and its twelve cuboctahedron
neighbors (degree 5). The update laws below act independently on each
connected component of the ball.

\subsection{Heat update and the monotonicity
lemma}\label{heat-update-and-the-monotonicity-lemma}

The heat (diffusion) update blends a field \(u_t \in \mathbb{R}^{N}\)
--- one value per lattice site, \(N\) the site count --- with its
symmetric-normalized adjacency average \(S = D^{-1/2} A D^{-1/2}\),
where \(A\) is the adjacency matrix of the move graph, \(D\) the
diagonal degree matrix, and rows of \(S\) are zero at isolated sites:

\begin{equation}
\label{eq:ivmdyn-heat}
u_{t+1} \;=\; (1-\alpha)\, u_t \;+\; \alpha\, S\, u_t ,
\qquad \alpha \in [0,1] .
\end{equation}

\textbf{Lemma (heat averaging is non-increasing in sum of squares).} For
every \(\alpha \in [0,1]\), every lattice built by
\passthrough{\lstinline!make\_lattice!}, and every state,

\begin{equation}
\label{eq:ivmdyn-lemma}
\lVert u_{t+1} \rVert_2^2 \;\le\; \lVert u_t \rVert_2^2 .
\end{equation}

\emph{Proof sketch.} \(S\) is symmetric with spectrum in \([-1,1]\) (it
is similar to the random-walk matrix \(D^{-1}A\)). The update operator
\((1-\alpha)I + \alpha S\) is therefore symmetric with eigenvalues in
\([1-2\alpha,\,1] \subseteq [-1,1]\), hence non-expansive in the 2-norm.
\(\square\)

This is exactly the claim the test suite asserts --- \emph{per step},
across seeds and the full coupling range
(\passthrough{\lstinline!test\_heat\_update\_l2\_nonincreasing\_per\_step!}
in \passthrough{\lstinline!tests/unit/lattice/test\_ivm\_dynamics.py!}).
Two scope restrictions are deliberate and tested:

\begin{enumerate}
\def\labelenumi{\arabic{enumi}.}
\tightlist
\item
  The guarantee is specific to the symmetric normalization \(S\). Plain
  row-average diffusion \(P = D^{-1}A\) on a truncated ball is
  \textbf{not} covered (boundary rows make \(P\) non-doubly-stochastic);
  the module neither uses nor claims it.
\item
  In the demo run (\(\alpha = 0.35\), \(T = 40\), seed 7) the sum of
  squares decreases monotonically from \(37.04\) to \(3.93\) while never
  increasing at any single step.
\end{enumerate}

\subsection{Majority update on integer
states}\label{majority-update-on-integer-states}

For integer-valued fields the rounded-averaging (``majority'') update is

\begin{equation}
\label{eq:ivmdyn-majority}
u_{t+1}(s) \;=\; \operatorname{rint}\!\Big( (1-\alpha)\, u_t(s)
\;+\; \alpha\, \tfrac{1}{\deg(s)} \!\!\sum_{n \in N(s)} u_t(n) \Big),
\end{equation}

with \(N(s)\) the neighbor set of site \(s\) and \(\deg(s) = |N(s)|\)
its degree, identity rows at isolated sites, and round-half-to-even ties
(\passthrough{\lstinline!numpy.rint!}), computed by
\passthrough{\lstinline!majority\_step!}. Every new value is a rounded
convex combination of the site and its neighbors, so the extremal
structure is preserved:

\textbf{Proposition (range containment).} Under
\eqref{eq:ivmdyn-majority} the field maximum is non-increasing and the
minimum non-decreasing, at every step, for every \(\alpha \in [0,1]\).

We deliberately make \textbf{no} sum-of-squares claim for this update:
rounding can increase it, and the test suite pins a demonstrating case
(origin at 0 with its twelve packers at 1 maps to thirteen sites at
value 1, raising the sum of squares from 12 to 13 ---
\passthrough{\lstinline!test\_majority\_step\_sum\_of\_squares\_can\_increase!}).
The honest summary: heat is \(L_2\)-monotone, majority is
range-preserving, and neither claim is stretched to cover the other
update.

\subsection{Learning the coupling from an observed
trajectory}\label{learning-the-coupling-from-an-observed-trajectory}

Given observed snapshots \(u_{\mathrm{obs}}(0..T)\) on a known lattice,
the coupling \(\alpha\) is identified gradient-free:
\passthrough{\lstinline!fit\_trajectory!} re-simulates the field from
the observed initial condition for each candidate coupling and minimizes
the trajectory mean squared error over the \(T\) post-initial snapshots
and the \(N\) sites, where \(u_{\alpha}(t)\) is the re-simulated field
under candidate coupling \(\alpha\):

\begin{equation}
\label{eq:ivmdyn-mse}
\mathcal{L}(\alpha) \;=\; \frac{1}{T\,N} \sum_{t=1}^{T}
\big\lVert u_{\alpha}(t) - u_{\mathrm{obs}}(t) \big\rVert_2^2 ,
\end{equation}

scanning the supplied grid, then re-gridding the interval between the
best candidate's grid neighbors \passthrough{\lstinline!refine\_rounds!}
times (grid/coordinate search; ties resolve to the smallest \(\alpha\);
fully deterministic). The result is a global optimum \emph{over the
evaluated candidates only} --- no convergence beyond the refinement
resolution is claimed. In the demo configuration (synthetic trajectory
generated at \(\alpha_{\text{true}} = 0.3\) on the \(R=3\) lattice,
\(T = 40\)), a 21-point grid plus refinement (33 MSE evaluations)
recovers \(\alpha \approx 0.300\) with
\(\mathcal{L} \approx 4.5 \times 10^{-33}\) --- machine precision for a
bitwise-deterministic re-simulation. When the true coupling lies
off-grid, refinement converges to within a grid-spacing fraction of it
(\passthrough{\lstinline!test\_fit\_refinement\_finds\_offgrid\_coupling!}:
0.37 recovered to \(\le 0.0125\) from a coarse
\(\{0, 0.25, 0.5, 0.75, 1\}\) grid).

\subsection{Demo figure}\label{demo-figure}

Generated by
\passthrough{\lstinline!uv run python quadmath/scripts/ivm\_dynamics\_demo.py!}
(\passthrough{\lstinline!MPLBACKEND=Agg!}); all randomness is seeded
(\passthrough{\lstinline!DynamicsParams.seed!}), so the figure is
reproducible byte-for-byte up to PNG encoding.

\begin{figure}
\centering
\pandocbounded{\includegraphics[keepaspectratio,alt={Dynamics and learning on the IVM lattice (R=3 ball, 13 sites). Top row: heat-field snapshots at t = 0, 20, 40 (\textbackslash alpha = 0.35, seed 7) over the embedded quadray sites (XYZ axes in embedding units); dot color encodes field value on a shared symmetric scale, and diffusion homogenizes the field by t = 40. Bottom left: sum-of-squares histories over t = 0 \textbackslash ldots 40 --- the heat curve decreases monotonically (lemma ); the majority curve falls steeply and plateaus, but unlike heat it carries no monotonicity guarantee, since rounding can increase the sum of squares (test\_majority\_step\_sum\_of\_squares\_can\_increase). Bottom middle: trajectory MSE versus coupling \textbackslash alpha on a logarithmic axis over the search grid, with the identified optimum (fit 0.3000) coinciding with the true generating \textbackslash alpha = 0.3 (dotted). Bottom right: the final integer majority field at t = 40. Generated by uv run python quadmath/scripts/ivm\_dynamics\_demo.py (MPLBACKEND=Agg, seeded via DynamicsParams.seed).}]{../output/figures/ivm_dynamics_demo.png}}
\caption{\textbf{Dynamics and learning on the IVM lattice (\(R=3\) ball,
13 sites)}. Top row: heat-field snapshots at \(t = 0, 20, 40\)
(\(\alpha = 0.35\), seed 7) over the embedded quadray sites (XYZ axes in
embedding units); dot color encodes field value on a shared symmetric
scale, and diffusion homogenizes the field by \(t = 40\). Bottom left:
sum-of-squares histories over \(t = 0 \ldots 40\) --- the heat curve
decreases monotonically (lemma \eqref{eq:ivmdyn-lemma}); the majority
curve falls steeply and plateaus, but unlike heat it carries no
monotonicity guarantee, since rounding can increase the sum of squares
(\protect\passthrough{\lstinline!test\_majority\_step\_sum\_of\_squares\_can\_increase!}).
Bottom middle: trajectory MSE versus coupling \(\alpha\) on a
logarithmic axis over the search grid, with the identified optimum
(\protect\passthrough{\lstinline!fit 0.3000!}) coinciding with the true
generating \(\alpha = 0.3\) (dotted). Bottom right: the final integer
majority field at \(t = 40\). Generated by
\protect\passthrough{\lstinline!uv run python quadmath/scripts/ivm\_dynamics\_demo.py!}
(\protect\passthrough{\lstinline!MPLBACKEND=Agg!}, seeded via
\protect\passthrough{\lstinline!DynamicsParams.seed!}).}
\end{figure}

\subsection{Reproducibility and test
contract}\label{reproducibility-and-test-contract}

\begin{itemize}
\tightlist
\item
  Tests:
  \passthrough{\lstinline!PYTEST\_DISABLE\_PLUGIN\_AUTOLOAD=1 uv run coverage run -m pytest tests/unit/lattice/test\_ivm\_dynamics.py -q!}
  then \passthrough{\lstinline!uv run coverage report!} --- 42 tests,
  100\% statement and branch coverage of
  \passthrough{\lstinline!src/quadmath/lattice/ivm\_dynamics.py!}, no
  mocks, all examples real numerics with fixed seeds.
\item
  Markdown:
  \passthrough{\lstinline!uv run python quadmath/scripts/validate\_markdown.py!}.
\item
  The two provable claims (lemma \eqref{eq:ivmdyn-lemma}, the range
  proposition under \eqref{eq:ivmdyn-majority}) are asserted per step;
  the non-claims (row-average diffusion, majority \(L_2\) behavior) are
  pinned by the honesty tests described above rather than left implicit.
\end{itemize}

\subsection{Cross-references}\label{cross-references}

\begin{itemize}
\tightlist
\item
  Site geometry and normalization: \href{03_quadray_methods.md}{Quadray
  Methods}.
\item
  Discrete point descent on the same move set:
  \href{04_optimization_in_4d.md}{Optimization in 4D}.
\item
  Field/energy vocabulary (free energy, Fisher information):
  \href{09_free_energy_active_inference.md}{Appendix B}; equation index:
  \href{08_equations_appendix.md}{Appendix A}.
\end{itemize}

\end{document}
